% problem-000207 10 阴离子传递多组分反应
\section*{10. 阴离子传递多组分反应}
利⽤阴离⼦传递(ARC)的多组分反应可以⾼效构筑复杂分⼦的⻣架体系。

\subsection*{10-1}
如下反应步骤是基于ARC的多组分反应（ASG为碳负离⼦稳定基团，E+为亲电试剂），写出此反应关键中间体M、N以及O的结构式。

（图：）

\subsection*{10-2}
根据以上信息，画出下列反应主要产物的结构式（要求⽴体化学）。

（图：）

\subsection*{10-3}
依据以上信息，画出构建如下产物所需三个反应原料的结构式（要求⽴体化学）。

（图：）

\subsection*{10-4}
将ARC策略和羟醛缩合反应相结合，可以⼀锅法⾼效构筑双环⻣架（如下图所⽰）。

（图：）

\subsection*{10-4}
-1 写出中间体P、Q及R的结构式（要求⽴体化学）。

\subsection*{10-4}
-2 判断S-V中的哪⼀个为此转换的主要产物。

V
（图：）

S

（图：）

T

（图：）

U

（图：）

\subsection*{10-5}
S和T两者之间的关系为：a) 对映异构体；b) 差向异构体；c) ⾮对映异构体。
